 nepředstavuje nezbytně hrozbu v~podobě plošného zániku pracovních míst, ale vyvolává hlubokou restrukturalizaci českého trhu práce. Dosavadní výzkumy i~modely potvrzují, že zatímco celková bilance počtu pracovních míst může zůstat relativně vyrovnaná, dopady na kvalifikační nároky, mzdovou stratifikaci a~především vstupní profesní pozice budou zásadní.~\cite{Weber_Industrie40_2015}~\cite{Massenkoff_LaborMarketAI_2026}\par

Technologický pokrok navíc vyvrací zažitou představu, že nejprve mizí pouze manuální rutinní práce na výrobních linkách. Zatímco průmyslová robotizace často naráží na mechanizační limity a~vysoké investiční náklady, generativní \acs{AI} s~minimálními marginálními náklady okamžitě nahrazuje kognitivní úkoly -- zákaznickou podporu, administrativní práce a~především juniorní kancelářské pozice, na nichž dříve mladá generace (např.~ze skupiny absolventů) získávala profesní zkušenosti a~kvalifikaci. Tento jev zásadně zvyšuje riziko stagnace celého segmentu mladých a~méně kvalifikovaných duševně pracujících, kteří bez silné státní podpory rekvalifikací a~zastoupení odborovými organizacemi ponesou veškeré náklady technologického přechodu.~\cite{Petit_AIEmploymentImpact_EU_2026}\par

Z~toho vyplývá nezbytnost, aby se budoucí kolektivní vyjednávání v~\acloc{ČR} neomezovalo pouze na tradiční vyjednávání o~odměňování. Kolektivní vyjednávaní musí zahrnovat pravidla pro zavádění systémů \acs{AI} na pracoviště, stanovení přísného dohledu nad intenzitou práce, právo na profesní rozšiřování kvalifikace hrazenou zaměstnavatelem a~spravedlivé rozdělení mimořádných výnosů ze zvýšené produktivity. Bez vytvoření moderního institucionálního rámce, který by tyto inovační aspekty zastřešil, hrozí, že technologická revoluce poukazované nerovnosti v~české ekonomice ještě více prohloubí.~\cite{OECD_GoodJobsForAll_2018}~\cite{NERV_NavrhyRustu2025}\par
